\documentclass[lettersize,journal]{IEEEtran}
\usepackage{amsmath,amsfonts}
\usepackage{array}
\usepackage{url}
\usepackage{graphicx}
\usepackage{booktabs}

\input{generated_metrics.tex}

\begin{document}

\title{Bias-Aware MFD-Based Macroscopic Traffic State Estimation and Latency-Compensated Perimeter Control under Non-Uniform Sparse Probes}

\author{Seyed Milad Jalilian
\thanks{The author is with the Department of Civil, Construction, Architectural and Environmental Engineering, University of Naples Federico II, Italy (e-mail: contact@miladjalilian.me). Simulation artifacts and manifests are released under the open-source \textit{PerimeterFlow} platform under Git commit \texttt{\ArtifactGitSHA} (status: \texttt{\ArtifactStatus}).}}

\markboth{IEEE Transactions on Intelligent Transportation Systems,~Technical Report}{Jalilian: Bias-Aware MFD Estimation and Perimeter Control}

\maketitle

\begin{abstract}
Feedback perimeter control utilizing the Macroscopic Fundamental Diagram (MFD) is a scalable method to mitigate metropolitan congestion. However, empirical deployment is fundamentally constrained by sparse probe penetration ($1\%$--$5\%$), topological corridor selection bias, GPS multipath noise, and cloud pipeline communication latency. This work formulates an integrated estimation and control architecture comprising: (1) a stateful kinematic quality gate, (2) a stratified link-class Horvitz-Thompson Edie estimator coupled with a two-branch hysteresis adaptive Kalman filter (PC-AKF), and (3) a residual-corrected Smith-PI perimeter controller. Evaluated across an $8\times 8$ two-reservoir SUMO grid across \EvalSeedCount{} held-out evaluation seeds and validated against empirical urban trajectory benchmarks from Amsterdam's municipal road network, the pipeline neutralizes corridor bias and maintains control loop stability.
\end{abstract}

\begin{IEEEkeywords}
Macroscopic Fundamental Diagram, Perimeter Control, Kalman Filter, Sparse Probes, Spatial Bias, Latency Compensation.
\end{IEEEkeywords}

\section{Introduction}
\IEEEPARstart{U}{rban} traffic networks experience non-linear capacity drop and gridlock when vehicular accumulation exceeds critical density \cite{geroliminis2008existence}. Perimeter control regulates flow across reservoir boundaries to maintain accumulation near optimal production thresholds \cite{keyvan2012exploiting}. Nevertheless, practical implementations rely on floating vehicle data exhibiting low penetration and severe spatial bias—clustering along arterial corridors while leaving local streets unmonitored \cite{ambuhl2016data}.

Furthermore, cloud message buffering and batch inference introduce telemetry latency $\tau_d \in [30\text{s}, 60\text{s}]$, inducing phase instability in standard proportional-integral (PI) controllers. This paper introduces \textit{PerimeterFlow}, an open-source research and simulation engine formulated to overcome these compounding constraints without requiring full inductive loop coverage.

\section{Methodology}
\subsection{Stateful Kinematic Quality Gate}
Raw probe pings $p_k = (t_k, x_k, y_k, v_k)$ are sequentially filtered against the last accepted valid anchor state $(\tilde{t}, \tilde{x}, \tilde{y}, \tilde{v})$ of that specific vehicle:
\begin{equation}
v_{\text{calc}} = \frac{\sqrt{(x_k - \tilde{x})^2 + (y_k - \tilde{y})^2}}{t_k - \tilde{t}}, \quad a_{\text{calc}} = \frac{v_{\text{calc}} - \tilde{v}}{t_k - \tilde{t}}
\end{equation}
Observations exceeding physical bounds ($v_{\text{calc}} > 22.22\text{ m/s}$, $a_{\text{calc}} \notin [-6.5, 4.5]\text{ m/s}^2$) are rejected without corrupting the anchor state, ensuring subsequent valid pings remain preserved.

\subsection{Stratified Horvitz-Thompson Estimator}
To correct for corridor over-representation, link classes $c \in \{\text{art}, \text{loc}\}$ with lane-length $L_c$ and observed length $L_c^{\text{obs}}(t)$ are weighted by inverse boundary-anchored class penetration $\hat{p}_c$:
\begin{equation}
\hat{N}(t) = \frac{1}{\Delta t} \sum_{c} \frac{L_c}{L_c^{\text{obs}}(t) \cdot \max(\hat{p}_c, p_{\min})} \sum_{i \in \mathcal{P}_c(t)} \tau_i^{\text{gated}}(t)
\end{equation}

\subsection{Hysteresis PC-AKF and Residual Smith-PI Control}
Vehicle accumulation is assimilated using a two-branch hysteresis production model $\Phi(N_t, \text{is\_unloading})$ \cite{saberi2014exploring}. Observation noise covariance is dynamically inflated using gate rejection rate $\rho_{\text{rej}}(t)$:
\begin{equation}
R_t = R_0 \left(\frac{p_{\text{ref}}}{\max(\hat{p}_{\text{eff}}(t), 0.005)}\right) \exp\left(\alpha_{\text{gate}} \rho_{\text{rej}}(t)\right)
\end{equation}
To compensate pipeline delay $\tau_d = d \cdot \Delta t$, a Smith predictor simulates forward rollout over past boundary actuations, applying model residual feedback $\delta_{t-d} = \hat{N}_{\text{AKF}}(t-d) - N_{\text{model}}(t-d)$ to govern normalized dimensionless gating control $u(t) \in [u_{\min}, 1.0]$.

\section{Experimental Protocol}
Simulations are conducted on an $8 \times 8$ two-reservoir SUMO grid with a 90-minute hysteresis demand profile. Per pre-registered plan \texttt{ANALYSIS\_PLAN.md}, calibration seeds 1--10 are utilized exclusively for parameter tuning, whereas evaluation seeds 11--40 ($n=\EvalSeedCount{}$) are held out for paired non-parametric inference. Comparisons report Hodges-Lehmann median difference ($\hat{\Delta}_{\text{HL}}$) and Cliff's delta ($\delta$) with Holm-Bonferroni correction.

\section{Results and Discussion}
\begin{table}[htbp]
\caption{Pre-Registered Evaluation Results Across $n=\EvalSeedCount{}$ Held-Out Seeds}
\label{tab:results}
\centering
\begin{tabular}{lcccc}
\toprule
Pipeline Variant & $\text{NRMSE}(N)$ & $\text{TND } [\text{veh}\cdot\text{h}]$ & Cliff's $\delta$ & Holm $p$ \\
\midrule
Unfiltered + Tuned PI & \MetricNRMSEBaseline & \MetricTNDBaseline & Baseline & --- \\
Proposed Smith-PI & \MetricNRMSESmith & \MetricTNDSmith & \MetricCliffsDelta & \MetricHolmP \\
\bottomrule
\end{tabular}
\end{table}

\section{Limitations and Threats to Validity}
While the framework preserves stability under latency, closed-loop gating cannot be counterfactually actuated on offline historical trajectory logs from Amsterdam's municipal network (NDW) \cite{knoop2013empirical}. Furthermore, boundary loop detectors remain necessary to anchor penetration rates $\hat{p}_c$.

\begin{thebibliography}{9}
\bibitem{geroliminis2008existence}
N. Geroliminis and C. F. Daganzo, ``Existence of urban-scale macroscopic fundamental diagrams: Some experimental findings,'' \emph{Transp. Res. Part B}, vol. 42, no. 9, pp. 759--770, 2008.

\bibitem{keyvan2012exploiting}
M. Keyvan-Ekbatani, A. Kouvelas, I. Papamichail, and M. Papageorgiou, ``Exploiting the fundamental diagram of urban networks for feedback-based gating,'' \emph{Transp. Res. Part B}, vol. 46, no. 10, pp. 1393--1403, 2012.

\bibitem{ambuhl2016data}
L. Amb{\"u}hl and M. Menendez, ``Data fusion algorithm for macroscopic fundamental diagram estimation,'' \emph{Transp. Res. Part C}, vol. 71, pp. 184--197, 2016.

\bibitem{saberi2014exploring}
M. Saberi and H. S. Mahmassani, ``Exploring properties of network-wide flow-density relations in a freeway network,'' \emph{Transp. Res. Rec.}, vol. 2315, pp. 153--163, 2012.

\bibitem{knoop2013empirical}
V.~L. Knoop and S.~P. Hoogendoorn, ``Empirical evaluation of the macroscopic fundamental diagram of a network,'' \emph{Transp. Res. Rec.}, vol. 2390, no. 1, pp. 62--70, 2013.
\end{thebibliography}

\end{document}
